\section{区域异质性}

\begin{frame}{异质性分析：减排效果主要体现在西部地区}
  \begin{table}
    \centering
    \small
    \begin{tabular}{@{}lccc@{}}
      \toprule
      地区 & 东部 & 中部 & 西部 \\
      \midrule
      绿色金融发展指数系数 & 不显著 & 不显著 & \textbf{-6.568*} \\
      减排效果判断 & 弱 & 弱 & 强 \\
      \bottomrule
    \end{tabular}
  \end{table}

  \begin{columns}[T,onlytextwidth]
    \column{0.48\textwidth}
    \begin{block}{西部地区为何更显著}
      \begin{itemize}
        \item 碳排放基数更高
        \item 边际减排空间更大
        \item 能源结构偏煤，改进空间明显
      \end{itemize}
    \end{block}

    \column{0.48\textwidth}
    \begin{exampleblock}{东中部为何不显著}
      \begin{itemize}
        \item 碳强度已相对较低
        \item 边际减排成本更高
        \item 绿色金融资源更多分散到非工业领域
      \end{itemize}
    \end{exampleblock}
  \end{columns}
\end{frame}

\begin{frame}{如何理解区域差异}
  \begin{itemize}
    \item 西部地区同时具备 \hlb{高排放基数 + 清洁能源项目集中布局} 两个条件。
    \item “西电东送”、大型风电光伏基地建设，使绿色金融资源更容易形成直接减排效果。
    \item 东中部地区产业结构更成熟，未来要释放绿色金融效能，关键在于 \hl{与产业升级、技术创新深度协同}。
  \end{itemize}

  \begin{alertblock}{研究含义}
    绿色金融政策有效，但其效果高度依赖地区的产业结构、能源禀赋与政策承接能力。
  \end{alertblock}
\end{frame}
